\DocumentMetadata{
  pdfversion=2.0,
  pdfstandard=ua-2,
  lang=en-US,
  tagging=on,
  tagging-setup={math/setup=mathml-SE,tabsorder=structure}
}
\documentclass[11pt,letterpaper]{article}
\usepackage[margin=0.68in,includefoot]{geometry}
\usepackage{newcomputermodern}
\usepackage{amsmath,amscd,mathtools}
\providecommand{\square}{\mdlgwhtsquare}
\usepackage[hidelinks]{hyperref}
\hypersetup{
  pdftitle={Algebra PhD Exam, May 2017},
  pdfauthor={University of Florida Department of Mathematics},
  pdfsubject={Graduate mathematics examination},
  pdfdisplaydoctitle=true,
  bookmarksopen=true
}
\setcounter{secnumdepth}{0}
\setlength{\parindent}{0pt}
\setlength{\parskip}{0.28em}
\setlength{\emergencystretch}{3em}
\setlength{\leftmargini}{2.15em}
\setlength{\leftmarginii}{2.65em}
\setlength{\leftmarginiii}{3em}
\renewcommand{\labelenumi}{\textbf{\theenumi.}}
\pagestyle{empty}
\raggedbottom
\allowdisplaybreaks
\newenvironment{examproblems}[1]{%
  \begin{enumerate}%
  \setcounter{enumi}{#1}%
  \setlength{\topsep}{0.35em}%
  \setlength{\partopsep}{0pt}%
  \setlength{\itemsep}{0.48em}%
  \setlength{\parsep}{0.18em}%
}{\end{enumerate}}
\newenvironment{examsubparts}{%
  \begin{enumerate}%
  \setlength{\topsep}{0.18em}%
  \setlength{\partopsep}{0pt}%
  \setlength{\itemsep}{0.16em}%
  \setlength{\parsep}{0.1em}%
}{\end{enumerate}}
\newenvironment{examinstructions}{%
  \par\begingroup\itshape
}{%
  \par\endgroup\vspace{0.18em}
}
\makeatletter
\renewcommand\section{\@startsection{section}{1}{0pt}%
  {-1.25ex plus -.25ex minus -.1ex}%
  {0.55ex plus .1ex}%
  {\normalfont\large\bfseries}}
\renewcommand{\@maketitle}{%
  \newpage\null\vskip -1.2em%
  {\footnotesize\bfseries
    \href{https://www.ufl.edu/}{University of Florida}, \href{https://math.ufl.edu/}{Department of Mathematics}\par}%
  \vskip 0.18em%
  \tagstructbegin{tag=Title}%
    {\LARGE\bfseries\href{https://gma.math.ufl.edu/exams/algebra-phd/}{Algebra PhD Exam}, May 2017\par}%
  \tagstructend%
  \vskip 0.35em%
  \tagmcbegin{artifact}%
    \hrule height 1pt%
  \tagmcend%
  \vskip 0.55em%
}
\makeatother
\title{Algebra PhD Exam, May 2017}
\author{}
\date{}
\begin{document}
\maketitle
\thispagestyle{empty}
\begin{examinstructions}
Answer seven problems. (If you turn in more, the first seven will be graded.) Within reason, you may quote theorems as long as you state them clearly.
\end{examinstructions}
\begin{examinstructions}
Note. Below ring means associative ring with identity, and module means unital module unless otherwise specified.
\end{examinstructions}
\begin{examproblems}{0}
\item
Let~\(F\) and~\(K\) be fields and let~\(S\) be a set of field homomorphisms from~\(F\) to~\(K\). Prove that~\(S\) is linearly independent over~\(K\) in the vector space of functions from~\(F\) to~\(K\).
\item
Let \(f(x)=x^3-10\).
\begin{examsubparts}
\item[\textnormal{(a)}]
Calculate the Galois group of \(f(x)\) over~\(\mathbb{Q}\).
\item[\textnormal{(b)}]
Calculate the Galois group of \(f(x)\) over \(\mathbb{Q}(\sqrt{2})\).
\end{examsubparts}
\item
Let \(\mathbb{Q}[x]\) be the algebra over~\(\mathbb{Q}\) of all polynomials in one variable over~\(\mathbb{Q}\), and similarly let \(\mathbb{Q}[x,y]\) be the algebra of polynomials in two variables. Prove that 
\[
\mathbb{Q}[x]\otimes_{\mathbb{Q}}\mathbb{Q}[x]\simeq\mathbb{Q}[x,y]
\]
 as algebras over~\(\mathbb{Q}\).
\item
A group~\(G\) is said to be nilpotent of class at most 2 if \([[G,G],G]=1\), where \([H,K]\) denotes the commutator group of the two subgroups~\(H\) and~\(K\) of~\(G\). Let~\(C\) be the concrete category whose objects are the nilpotent groups of class at most 2 and whose morphisms are all the group homomorphisms between them.
\begin{examsubparts}
\item[\textnormal{(a)}]
Prove that, for every set~\(S\), there is a free object in~\(C\) with free generator set~\(S\).
\item[\textnormal{(b)}]
Prove that if \(S\ne\varnothing\), then the free object in~\(C\) with free generator set~\(S\) is infinite.
\end{examsubparts}
\item
Let~\(R\) be a ring. Prove that the following two conditions are equivalent for a left~\(R\)-module~\(P\). (You may not assume any properties of projective modules.)
\begin{examsubparts}
\item[\textnormal{(a)}]
Every short exact sequence of~\(R\)-modules 
\[
0\to A\to B\to P\to 0
\]
 splits;
\item[\textnormal{(b)}]
There is a free~\(R\)-module~\(F\) and an~\(R\)-module~\(K\) such that \(F\simeq K\oplus P\).
\end{examsubparts}
\item
Let~\(R\) be the ring of all rational numbers with odd denominators. Prove that the Jacobson radical \(J(R)\) of~\(R\) consists of all fractions whose denominator is odd and whose numerator is even.
\item
Let~\(R\) be an integral domain, and let~\(I\) be an invertible ideal of~\(R\). Prove that~\(I\) is finitely generated.
\item
Let~\(R\) be a commutative ring with identity, let~\(M\) be a maximal ideal of~\(R\), and let~\(n\) be a positive integer. Prove \(R/M^n\) has a unique prime ideal, so that, in particular, \(R/M^n\) is a local ring.
\item
Prove the following version of Nakayama’s Lemma. Let~\(A\) be a commutative ring with identity, let~\(M\) be a finitely generated~\(A\)-module, and let~\(I\) be an ideal contained in the intersection of all the maximal ideals of~\(A\). Show that if \(IM=M\) then \(M=\{0\}\).
\item
Let~\(R\) be a ring (possibly non-commutative, and possibly without identity). Let~\(M\) be an irreducible left~\(R\)-module. Let \(D=\operatorname{End}_R(M)\).
\begin{examsubparts}
\item[\textnormal{(a)}]
Prove \(D=\operatorname{End}_R(M)\) is a division ring. (You may NOT assume Schur’s Lemma for this.)
\item[\textnormal{(b)}]
Give an example of a ring~\(R\) and an irreducible module~\(M\) such that~\(D\) is not commutative.
\end{examsubparts}
\item
Classify (up to ring isomorphism) all semisimple rings of order \(720\).
\end{examproblems}
\end{document}
